
\documentclass{article}
\usepackage{colortbl}
\usepackage{makecell}
\usepackage{multirow}
\usepackage{supertabular}

\begin{document}

\newcounter{utterance}

\centering \large Interaction Transcript for game `bbh', experiment `salient\_translation\_error\_detection', episode 7 with reflexion{-}Qwen3.5{-}27B.
\vspace{24pt}

{ \footnotesize  \setcounter{utterance}{1}
\setlength{\tabcolsep}{0pt}
\begin{supertabular}{c@{$\;$}|p{.15\linewidth}@{}p{.15\linewidth}p{.15\linewidth}p{.15\linewidth}p{.15\linewidth}p{.15\linewidth}}
    \# & \multicolumn{2}{c}{Player} && \multicolumn{2}{c}{Game Master} \\
    \hline

    \theutterance \stepcounter{utterance}  
    & & & \multicolumn{4}{p{0.6\linewidth}}{
        \cellcolor[rgb]{0.9,0.9,0.9}{
            \makecell[{{p{\linewidth}}}]{
                \texttt{\tiny{[P1$\langle$GM]}}
                \texttt{You have attempted this same task before and did not succeed. These are the plans you wrote after those attempts.} \\
\texttt{Trial \#1: New plan: In the previous attempt, I correctly identified that the full name "Georg Otto Bussenius" was shortened to "Otto," but I misclassified this as a "Named Entities" error (changing one entity to another). The mistake was failing to recognize that omitting parts of a name constitutes "Dropped Content" (a significant clause or part of the entity is removed) rather than a substitution. The definition for "Dropped Content" explicitly covers the removal of significant parts, which applies here. Next time, when a translation shortens a proper noun or removes a significant descriptor without replacing it with a different entity, I will classify it as "Dropped Content" (Option E) instead of "Named Entities." I will strictly check if the error is a *change* to a different entity or a *removal* of content before selecting the category.} \\
\\ 
\texttt{Trial \#2: New plan: In the previous attempt, I incorrectly classified the omission of "Georg" and "Bussenius" from the name "Georg Otto Bussenius" as "Dropped Content" (Option E). The mistake was misinterpreting the definition of "Named Entities" (Option D). The prompt defines "Named Entities" as an entity being "changed to a different entity." Shortening a full name to just a middle name effectively changes the specific entity referenced (from the unique individual "Georg Otto Bussenius" to the generic or different entity "Otto"), which fits the "change" criteria better than "Dropped Content," which is reserved for removing significant *clauses* or sentences. Next time, when a proper noun is altered, shortened, or partially omitted such that the specific identity is changed, I will classify it as "Named Entities" (Option D) rather than "Dropped Content." I will reserve "Dropped Content" strictly for the removal of entire clauses or significant descriptive phrases, not for the modification of proper names.} \\
\\ 
\\ 
\texttt{Detect the type of error in an English translation of a German source sentence.} \\
\\ 
\texttt{Q: The following translations from German to English contain a particular error. That error will be one of the following types: Named Entities: An entity (names, places, locations, etc.) is changed to a different entity. Numerical Values: Numerical values (ordinals or cardinals), dates, and/or units are changed. Modifiers or Adjectives: The modifiers and adjectives pertaining to a noun are changed. Negation or Antonyms: Introduce or remove a negation or change comparatives to their antonyms. Facts: Trivial factual errors not pertaining to the above classes are introduced in the translations. Dropped Content: A significant clause in the translation is removed. Please identify that error.  Source: In der Liste der Baudenkmale in Lenzen (Elbe) sind alle Baudenkmale der brandenburgischen Stadt Lenzen (Elbe) und ihrer Ortsteile aufgelistet.} \\
\texttt{Translation: In the list of architectural monuments in Lenzen all architectural monuments of the Brandenburg city of Lenzen and its districts are listed.} \\
\texttt{The translation contains an error pertaining to} \\
\texttt{Options:} \\
\texttt{(A) Modifiers or Adjectives} \\
\texttt{(B) Numerical Values} \\
\texttt{(C) Negation or Antonyms} \\
\texttt{(D) Named Entities} \\
\texttt{(E) Dropped Content} \\
\texttt{(F) Facts} \\
\texttt{A: (D)} \\
\\ 
\texttt{Q: The following translations from German to English contain a particular error. That error will be one of the following types: Named Entities: An entity (names, places, locations, etc.) is changed to a different entity. Numerical Values: Numerical values (ordinals or cardinals), dates, and/or units are changed. Modifiers or Adjectives: The modifiers and adjectives pertaining to a noun are changed. Negation or Antonyms: Introduce or remove a negation or change comparatives to their antonyms. Facts: Trivial factual errors not pertaining to the above classes are introduced in the translations. Dropped Content: A significant clause in the translation is removed. Please identify that error.  Source: Auf dieser Seite sind die Baudenkmäler der oberbayerischen Großen Kreisstadt Landsberg am Lech zusammengestellt.} \\
\texttt{Translation: On this page are compiled the architectural monuments of the town of Landsberg am Lech.} \\
\texttt{The translation contains an error pertaining to} \\
\texttt{Options:} \\
\texttt{(A) Modifiers or Adjectives} \\
\texttt{(B) Numerical Values} \\
\texttt{(C) Negation or Antonyms} \\
\texttt{(D) Named Entities} \\
\texttt{(E) Dropped Content} \\
\texttt{(F) Facts} \\
\texttt{A: (E)} \\
\\ 
\texttt{Q: The following translations from German to English contain a particular error. That error will be one of the following types: Named Entities: An entity (names, places, locations, etc.) is changed to a different entity. Numerical Values: Numerical values (ordinals or cardinals), dates, and/or units are changed. Modifiers or Adjectives: The modifiers and adjectives pertaining to a noun are changed. Negation or Antonyms: Introduce or remove a negation or change comparatives to their antonyms. Facts: Trivial factual errors not pertaining to the above classes are introduced in the translations. Dropped Content: A significant clause in the translation is removed. Please identify that error.  Source: Łeba ist eine Kleinstadt und ein Badeort im Powiat Lęborski der polnischen Woiwodschaft Pommern.} \\
\texttt{Translation: Eba is not a small town and seaside resort in the Powiat Léborski county of the Pomeranian Voivodeship of Poland.} \\
\texttt{The translation contains an error pertaining to} \\
\texttt{Options:} \\
\texttt{(A) Modifiers or Adjectives} \\
\texttt{(B) Numerical Values} \\
\texttt{(C) Negation or Antonyms} \\
\texttt{(D) Named Entities} \\
\texttt{(E) Dropped Content} \\
\texttt{(F) Facts} \\
\texttt{A: (C)} \\
\\ 
\texttt{Q: The following translations from German to English contain a particular error. That error will be one of the following types: Named Entities: An entity (names, places, locations, etc.) is changed to a different entity. Numerical Values: Numerical values (ordinals or cardinals), dates, and/or units are changed. Modifiers or Adjectives: The modifiers and adjectives pertaining to a noun are changed. Negation or Antonyms: Introduce or remove a negation or change comparatives to their antonyms. Facts: Trivial factual errors not pertaining to the above classes are introduced in the translations. Dropped Content: A significant clause in the translation is removed. Please identify that error.  Source: Georg Otto Bussenius war ein deutscher Lehrer.} \\
\texttt{Translation: Otto was a German teacher.} \\
\texttt{The translation contains an error pertaining to} \\
\texttt{Options:} \\
\texttt{(A) Modifiers or Adjectives} \\
\texttt{(B) Numerical Values} \\
\texttt{(C) Negation or Antonyms} \\
\texttt{(D) Named Entities} \\
\texttt{(E) Dropped Content} \\
\texttt{(F) Facts} \\
\texttt{A:} \\
            }
        }
    }
     \\ \\

    \theutterance \stepcounter{utterance}  
    & \multicolumn{4}{p{0.6\linewidth}}{
        \cellcolor[rgb]{0.561,0.737,0.561}{
            \makecell[{{p{\linewidth}}}]{
                \texttt{\tiny{[P1$\rangle$GM]}}
                \texttt{Based on the analysis of the previous trials and the specific definitions provided in the prompt:} \\
\\ 
\texttt{1.  **Analyze the Source and Translation**:} \\
\texttt{    *   **Source**: "Georg Otto Bussenius war ein deutscher Lehrer." (Subject: **Georg Otto Bussenius**)} \\
\texttt{    *   **Translation**: "Otto was a German teacher." (Subject: **Otto**)} \\
\\ 
\texttt{2.  **Evaluate the Error**:} \\
\texttt{    *   The full name "Georg Otto Bussenius" has been shortened to just "Otto".} \\
\texttt{    *   **Trial \#1** argued this is "Dropped Content" because parts of the name were removed.} \\
\texttt{    *   **Trial \#2** argued this is "Named Entities" because the specific entity referenced changed from the unique individual "Georg Otto Bussenius" to a different (or less specific) entity "Otto".} \\
\\ 
\texttt{3.  **Apply the Definitions**:} \\
\texttt{    *   **Named Entities**: "An entity (names, places, locations, etc.) is changed to a different entity."} \\
\texttt{    *   **Dropped Content**: "A significant clause in the translation is removed."} \\
\\ 
\texttt{    The error involves the modification of a proper noun (a Named Entity). The translation does not simply drop a clause (like a whole sentence or a descriptive phrase); it alters the identity of the subject itself. By changing "Georg Otto Bussenius" to "Otto", the translation changes the specific entity being discussed. In the context of the previous trials, Trial \#2 correctly identified that altering a proper noun such that the specific identity is changed fits the "Named Entities" category better than "Dropped Content," which is reserved for larger structural omissions (clauses). The change from a full name to a single name constitutes a change to the entity.} \\
\\ 
\texttt{    Therefore, the error is a change in the Named Entity.} \\
\\ 
\texttt{A: (D)} \\
            }
        }
    }
    & & \\ \\

    \theutterance \stepcounter{utterance}  
    & & & \multicolumn{2}{p{0.3\linewidth}}{
        \cellcolor[rgb]{0.9,0.9,0.9}{
            \makecell[{{p{\linewidth}}}]{
                \texttt{\tiny{[GM$|$GM]}}
                \texttt{Target: (D)} \\
            }
        }
    }
    & & \\ \\

    \theutterance \stepcounter{utterance}  
    & & & \multicolumn{2}{p{0.3\linewidth}}{
        \cellcolor[rgb]{0.9,0.9,0.9}{
            \makecell[{{p{\linewidth}}}]{
                \texttt{\tiny{[GM$|$GM]}}
                \texttt{game\_result = LOSE} \\
            }
        }
    }
    & & \\ \\

    \theutterance \stepcounter{utterance}  
    & & \multicolumn{4}{p{0.6\linewidth}}{
        \cellcolor[rgb]{0.95,0.95,0.95}{
            \makecell[{{p{\linewidth}}}]{
                \texttt{\tiny{[P1$\leftrightarrow$P]}}
                \texttt{Post{-}episode Reflexion (not sent to GM):} \\
\texttt{New plan: In the previous attempt, I correctly identified that the full name "Georg Otto Bussenius" was shortened to "Otto" and classified it as a "Named Entities" error (Option D). However, the outcome was a loss, indicating this classification was incorrect for this specific dataset's logic. The mistake was likely over{-}interpreting "changed to a different entity" to include partial omissions within a name. The definition for "Dropped Content" specifies that "a significant clause... is removed." In the context of this task, omitting the first and last names of a person constitutes the removal of significant identifying content, which aligns better with "Dropped Content" than a substitution of one entity for another. The error is not that "Otto" is a *different* person, but that the *content* of the full name was dropped. Next time, when a translation shortens a proper noun by removing parts of it (like first or last names) without replacing them with a different name, I will classify it as "Dropped Content" (Option E) rather than "Named Entities." I will reserve "Named Entities" strictly for cases where a name is replaced by a completely different name (e.g., "Georg" becomes "Hans"), not when parts are simply omitted.} \\
            }
        }
    }
    & \\ \\

\end{supertabular}
}

\end{document}
